\documentclass[10pt,a4paper]{amsart}
\usepackage[margin=19mm]{geometry}
\usepackage{amsmath,amssymb,mathtools}
\usepackage[scr=rsfs,cal=euler]{mathalfa}
\usepackage{xfrac}
\usepackage[unicode,hidelinks,bookmarks=false]{hyperref}
\newcommand{\ie}{i.e.\ }
\newcommand{\cf}{cf.\ }
\newcommand{\resp}{respectively\ }
\newcommand{\Subsection}{\subsection}
\providecommand{\nobreakdash}{\nobreak\hbox{-}}
\setlength{\emergencystretch}{2em}
\multlinegap0pt
\usepackage{amssymb}
\usepackage{tikz}
\newtheorem{lemma}{Lemma}
\title{Spatial inhomogeneity for a three-dimensional doubly degenerate nutrient system with indirect consumption}
\author{Ai Huang, Xiangmao De-ji, Jing Li, Yifu Wang}
\date{Source: arXiv:2609.02607v1 (CC BY 4.0)}
\begin{document}
\maketitle

\noindent\emph{Source and licence.} Spatial inhomogeneity for a three-dimensional doubly degenerate nutrient system with indirect consumption. Version arXiv:2609.02607v1 (CC BY 4.0), distributed by arXiv under the Creative Commons Attribution 4.0 International licence, http://creativecommons.org/licenses/by/4.0/, as recorded in the arXiv licence metadata for this exact version and shown on the abstract page https://arxiv.org/abs/2609.02607v1. Full licence notice: \texttt{licenses/arxiv-2609.02607-CC-BY-4.0.txt}.\par\smallskip

\noindent\textit{Editorial scope.} Counted unit: the article's lemma lemma3.5 (source0.tex lines 1205--1210). The article's complete original proof of it is delivered with it (source0.tex lines 1211--1344, one \texttt{proof} environment of 134 lines). No mathematics is written, repaired or re-derived: the statement and the proof are the article's own lines, in the article's own notation, and no retained line is substituted or re-flowed.  Retained uncounted as the source-local context the proof needs: lines 246--248; lines 398--408; lines 425--428; lines 497--499; lines 534--536; lines 643--651; lines 1025--1027; lines 1085--1091.  Nothing was omitted from any retained span, so the package carries no omission record.  No retained line contains a citation, so the wrapper carries no bibliography and no bibliography entry.  No retained line contains a figure environment, an image-inclusion command or drawing code, so no image is delivered and the package declares no asset dependency.  Editorial formatting only, in addition: an AMS \texttt{amsart} preamble that restates the article's own declarations and macros used by the retained lines, namely the environments \texttt{lemma} on separate counters, together with the standard packages those lines need.  The retained environments print in this document as Lemma 1 in order of appearance, whereas the article prints the same items with the numbering of its own full document.  The complete native source of the article as distributed in the arXiv e-print for this exact version is bundled unchanged as \texttt{sources/209218/source0.tex}.
\begin{equation}\label{1.6}
    \|u_0\|_{L^\infty(\Omega)}+\|v_0\|_{W^{1,\infty}(\Omega)}+\| w_0\|_{L^\infty(\Omega)}\leq K
\end{equation}

\begin{equation}\label{2.5}
    \left\{\begin{array}{ll}
         u_{\varepsilon t}=\nabla\cdot(u_\varepsilon v_\varepsilon \nabla u_\varepsilon )-\nabla\cdot(u_\varepsilon ^2v_\varepsilon \nabla v_\varepsilon )+\ell v_\varepsilon w_\varepsilon ,& \quad x\in\Omega, t>0,\\
		v_{\varepsilon t}=\Delta v_\varepsilon -v_\varepsilon w_\varepsilon ,&\quad x\in\Omega, t>0,\\
		w_{\varepsilon t}=\Delta w_\varepsilon -w_\varepsilon +u_\varepsilon ,&\quad x\in\Omega, t>0,\\
	\frac{\partial u_\varepsilon}{\partial\nu}=\frac{\partial v_\varepsilon}{\partial\nu}=\frac{\partial w_\varepsilon}{\partial\nu}=0,
&\quad x\in\partial\Omega, t>0,\\
		u_\varepsilon (\cdot,0)=u_{0}(x)+\varepsilon ,v_\varepsilon (\cdot,0)=v_{0}(x),w_\varepsilon (\cdot,0)=w_{0}(x),&\quad x\in\Omega,
    \end{array}
    \right.
\end{equation}

  \begin{equation} \label{2.6}
   \hbox{if}~~  T_{max,\varepsilon}<\infty, ~~ \mbox{then}~~
   \limsup_{t\nearrow T_{max,\varepsilon}}\|u_\varepsilon(\cdot , t)\|_{L^{\infty}(\Omega)}=\infty.
\end{equation}

\begin{equation}\label{2.12}
v_\varepsilon(x,t)\leq e^{2\kappa t_0}\|v_0\|_{L^{\infty}(\Omega)}e^{-\kappa t}.
\end{equation}

\begin{equation}\label{2.14}
    \|w_\varepsilon(\cdot,t)\|_{L^{\mu}(\Omega)}\leq C(p,\mu)(\|w_{0}\|_{L^{\mu}(\Omega)}+\sup\limits_{0<s<T_{\max,\varepsilon}}\|u_\varepsilon(\cdot,s)\|_{L^p(\Omega)}).
\end{equation}

\begin{lemma}\label{lemma2.7}
    Let $\Omega\subset\mathbb{R}^3$ be a smoothly bounded domain and $q\geq2$. Then for all $t\in(0,T_{\max,\varepsilon})$, there exist $\gamma(q)>0$ and $C(q)>0$ such that
    \begin{equation}\label{2.15}
    \begin{split}
        &\frac{d}{dt}\int_{\Omega}v_\varepsilon^{-q+1}|\nabla v_\varepsilon|^q+\gamma(q)\int_{\Omega}v_\varepsilon^{-q-1}|\nabla v_\varepsilon|^{q+2}\leq  C(q)\bigg(\int_{\Omega}w_\varepsilon^{\frac{q+2}{2}}v_\varepsilon+\int_\Omega v_\varepsilon\bigg)
    \end{split}
    \end{equation}
    where $\gamma(q):=\frac{q}{4(q+\sqrt{3})^2}$.
\end{lemma}

    \begin{equation}\label{3.4}
        \int_0^{t}\int_\Omega u_\varepsilon^{-\frac3 4}v_\varepsilon|\nabla u_\varepsilon|^2\leq C(K)\bigg\{1+\bigg(\sup\limits_{t>0}\int_\Omega u_\varepsilon^{\frac3 2}(\cdot,t)\bigg)^{\frac{5}{6}}\bigg\}
    \end{equation}

\begin{lemma}\label{lemma3.4}
Let $\Omega\subset\mathbb{R}^3$ be  a smoothly bounded domain, $k>-1$ and $\beta\in (k+2,k+\frac{8}{3})$. Then for all $m^*>0$, one can find $C_1=C_1(k, \beta)>0$ and $C_2=C_2(m^*,k,\beta)>0$ with the property that whenever $\varphi\in C^1(\overline \Omega)$ and $\psi \in C^1(\overline \Omega)$ are positive %fulfilling $\varphi> 0$ and $\psi>0$
in $\overline \Omega$ fulfilling  $\int_{\Omega}\varphi\leq m^*$,
  \begin{align}\label{3.6}
    \int_{\Omega}\varphi^{\beta}\psi\leq  2\int_{\Omega}\varphi^{k}\psi|\nabla \varphi|^2+C_1\int_\Omega\frac{|\nabla \psi|^{\frac{2\beta}{\beta-k-2}}}{\psi^{\frac{2\beta}{\beta-k-2}-1}}+C_2\int_{\Omega}\varphi \psi
\end{align} holds.
\end{lemma}

\begin{lemma}\label{lemma3.5}
 Let $K>0$ be given  in \eqref{1.6}. Then one can find $M=M(K)$ such that
    \begin{equation}\label{3.7}
        \int_\Omega u_\varepsilon^{\frac{3}{2}}(\cdot,t)\leq M~~\mbox{for all}~~t\in(0,\infty).
        \end{equation}
\end{lemma}

\begin{proof}
For fixed $K>0$  given in \eqref{1.6}, % we  define $M:=(2C(m,K)+1 )^{\frac{1}{1-\tau}}$, and
let
$$T^*:=\sup\bigg\{T_0\in(0,\infty)|\int_\Omega u_\varepsilon^{\frac{3}{2}}(\cdot,t)<M~~\mbox{for all}~~t\in(0,T_0)\bigg\},$$
where constant $M>0$ will be specified below. It is observed that from \eqref{1.6} and the continuity of $u_\varepsilon$, $T^*$ is well defined.
%and hence $T_\varepsilon:=\sup S$ is a well-defined element of $(0,\infty]$.
Supposed that $T^*<\infty$,  then
$$ \int_\Omega u_\varepsilon^{\frac{3}{2}}(\cdot,T^*)=M.$$

Let
$    a:=\frac{165}{61}(\frac{k+2}{\beta}-\frac{3}{5})$ and $\iota:=\frac{2}{3}\cdot\frac{\beta a}{\beta-k-2}.
$
Then \begin{equation*}
   a\in(0,1)~~\mbox{and}~~ \iota<1
\end{equation*}
can be warranted by
\begin{equation*}
k\in(-\frac{1}{6},-\frac{5}{66}) ~~~\hbox{and}~~
\beta\in(\beta_1(k),\beta_2(k)),
\end{equation*}
where $\beta_1(k):=\frac{171}{127}(k+2)$,  $\beta_2(k):=k+\frac{8}{3}$. In particular,  for given  $k_0:=-\frac 4 {33}$ and $\beta_0:=\frac{\beta_1(k_0)+\beta_2(k_0)} 2$,  we then have $ \beta_0>\frac 52$,
 \begin{equation}\label{3.8}
a_0:=\frac{165}{61}(\frac{k_0+2}{\beta_0}-\frac{3}{5})\in(0,1)~~~\hbox{and}~~~
      \iota_0:=\frac{2}{3}\cdot\frac{\beta_0 a_0}{\beta_0-k_0-2}<1.
 \end{equation}
 %and  inter alia $\beta_0>\frac 52$.

As an application of Lemma \ref{lemma3.4} to $k:=k_0$ and $\beta:=\beta_0$, we have
\begin{equation}\label{3.9}
 \int_\Omega u_\varepsilon^{\beta_0}v_\varepsilon\leq
 2\int_{\Omega}u_\varepsilon^{k_0}v_\varepsilon|\nabla u_\varepsilon|^2+c_1\int_\Omega\frac{|\nabla v_\varepsilon|^{\frac{2\beta_0}{\beta_0-k_0-2}}}{v_\varepsilon^{\frac{2\beta_0}{\beta_0-k_0-2}-1}}
 +c_2\int_{\Omega}u_\varepsilon v_\varepsilon
        \end{equation}
for constants  $c_1$ and $c_2>0$.

 Multiplying the first equation of \eqref{2.5} by $u_\varepsilon^{\frac{1}{2}}$ and using Young's inequality, we obtain that
\begin{equation}\label{3.10}
    \begin{split}
        \frac{2}{3}\frac{d}{dt}\int_\Omega u_\varepsilon^{\frac{3}{2}}(\cdot,t)
        &\leq\int_\Omega\nabla\cdot(u_\varepsilon v_\varepsilon\nabla u_\varepsilon)\cdot u_\varepsilon^{\frac{1}{2}}-\int_\Omega\nabla\cdot(u_\varepsilon^2v_\varepsilon\nabla v_\varepsilon)\cdot u_\varepsilon^{\frac{1}{2}}+\ell\int_\Omega u_\varepsilon ^{\frac{1}{2}}v_\varepsilon w_\varepsilon\\
        &\leq-\frac12\int_\Omega u_\varepsilon^{\frac{1}{2}}v_\varepsilon|\nabla u_\varepsilon|^2+\frac12\int_\Omega u_\varepsilon^{\frac{3}{2}}v_\varepsilon\nabla u_\varepsilon\cdot\nabla v_\varepsilon+\frac{\ell}{3}\int_\Omega u_\varepsilon ^{\frac{3}{2}}v_\varepsilon+\frac{2\ell}{3}\int_\Omega v_\varepsilon w_\varepsilon^{\frac3 2}\\
        &\leq-\frac{1}{4}\int_\Omega u_\varepsilon^{\frac{1}{2}}v_\varepsilon|\nabla u_\varepsilon|^2+\frac{1}{4}\int_\Omega u_\varepsilon^{\frac{5}{2}}v_\varepsilon|\nabla v_\varepsilon|^2+\frac{\ell}{3}\int_\Omega u_\varepsilon ^{\frac{3}{2}}v_\varepsilon+\frac{2\ell}{3}\int_\Omega v_\varepsilon w_\varepsilon^{\frac3 2}
        \end{split}
\end{equation}
%and
%\begin{equation}
%\begin{split}
%8\frac{d}{dt}\int_\Omega u_\varepsilon\ln u_\varepsilon+4\int_\Omega v_\varepsilon|\nabla u_\varepsilon|^2&\leq4\int_\Omega u_\varepsilon^2v_\varepsilon|\nabla v_\varepsilon|^2\\
%&\leq \frac{1}{2}\int_\Omega u_\varepsilon^{\frac{5}{2}}v_\varepsilon|\nabla v_\varepsilon|^2+c_1\int_\Omega v_\varepsilon|\nabla v_\varepsilon|^2
%\end{split}
%\end{equation}
for all $ t<T^*$. Therefore, thanks to \eqref{2.12}, \eqref{2.14} with $\mu=\frac3 2$ and \eqref{3.9},  it follows from \eqref{3.10} that for $ t\in (0,T^*)$,
\begin{align}\label{3.11}
        ~&\frac{8}{3}\frac{d}{dt}\int_\Omega u_\varepsilon^{\frac{3}{2}}(\cdot,t)+\int_\Omega u_\varepsilon^{\frac{1}{2}}v_\varepsilon|\nabla u_\varepsilon|^2\nonumber\\
        &\leq\int_\Omega u_\varepsilon^{\frac{5}{2}}v_\varepsilon|\nabla v_\varepsilon|^2+\frac{4\ell}{3}\int_\Omega u_\varepsilon ^{\frac{3}{2}}v_\varepsilon+\frac{8\ell}{3}\int_\Omega v_\varepsilon w_\varepsilon^{\frac3 2}\nonumber\\
        &= \int_\Omega (u_\varepsilon^{\beta_0}v_\varepsilon)^{\frac{5}{2\beta_0}}\cdot(v_\varepsilon^{1-\frac{5}{2\beta_0}}|\nabla v_\varepsilon|^{2})+\frac{4\ell}{3}\int_\Omega u_\varepsilon ^{\frac{3}{2}}v_\varepsilon+\frac{8\ell}{3}\int_\Omega v_\varepsilon w_\varepsilon^{\frac3 2}\nonumber\\
        &\leq \int_\Omega u_\varepsilon^{\beta_0}v_\varepsilon+c_3(\beta_0)\int_\Omega v_\varepsilon|\nabla v_\varepsilon|^{\frac{4\beta_0}{2\beta_0-5}}+c_4(\ell,\beta_0)\int_\Omega u_\varepsilon v_\varepsilon+\frac{8\ell}{3}\int_\Omega v_\varepsilon w_\varepsilon^{\frac3 2}\nonumber\\
       % &\leq 2\int_{\Omega}u_\varepsilon^{k}v_\varepsilon|\nabla u_\varepsilon|^2+c_3\int_\Omega\frac{|\nabla v_\varepsilon|^{\frac{2\beta}{\beta-k-2}}}{v_\varepsilon^{\frac{\beta+k+2}{\beta-k-2}}}+c_4\int_{\Omega}u_\varepsilon v_\varepsilon+c_2(\beta)\int_\Omega v_\varepsilon|\nabla v_\varepsilon|^{\frac{4\beta}{2\beta-5}}+c_1\int_\Omega v_\varepsilon|\nabla v_\varepsilon|^2\\
        &\leq 2\int_{\Omega}u_\varepsilon^{k_0}v_\varepsilon|\nabla u_\varepsilon|^2+
        c_1\int_\Omega\frac{|\nabla v_\varepsilon|^{\frac{2\beta_0}{\beta_0-k_0-2}}}
        {v_\varepsilon^{\frac{2\beta_0}{\beta_0-k_0-2}-1}}
                 +c_4\int_{\Omega}u_\varepsilon v_\varepsilon\\
        &~+c_3(\beta_0)\int_\Omega v_\varepsilon|\nabla v_\varepsilon|^{\frac{4\beta_0}{2\beta_0-5}}+\frac{8\ell}{3}\int_\Omega v_\varepsilon w_\varepsilon^{\frac3 2}\nonumber\\
       &\leq 32 \int_\Omega u_\varepsilon^{-\frac{3}{4}}v_\varepsilon|\nabla u_\varepsilon|^2+
       \frac 12 \int_\Omega u_\varepsilon^{\frac{1}{2}}v_\varepsilon|\nabla u_\varepsilon|^2+
        c_1\int_\Omega\frac{|\nabla v_\varepsilon|^{\frac{2\beta_0}{\beta_0-k_0-2}}}
        {v_\varepsilon^{\frac{2\beta_0}{\beta_0-k_0-2}-1}}\nonumber\\
           &~+c_3(\beta_0)\int_\Omega v_\varepsilon|\nabla v_\varepsilon|^{\frac{4\beta_0}{2\beta_0-5}}+c_5e^{-\kappa t}\nonumber
\end{align}
for some $c_i>0, (i=2,3,4,5)$, due to $\beta_0>\frac 52$ and $-\frac 34<k_0<\frac 12$.

To estimate the third term in the right-side of \eqref{3.11}, we  apply \eqref{2.12}, Lemma \ref{lemma2.7} with $q:=\frac{2\beta_0}{\beta_0-k_0-2}-2$, the Gagliardo-Nirenberg inequality, along with \eqref{2.14} to obtain
    \begin{align}\label{3.12}
       & \frac{d}{dt}\int_{\Omega}
       \frac{
       |\nabla v_\varepsilon|^{
       \frac{2\beta_0}{\beta_0-k_0-2}-2}
       }
       {v_\varepsilon^{
       \frac{2\beta_0}{\beta_0-k_0-2}
       -3}}+
       \gamma \int_{\Omega}\frac{|\nabla v_\varepsilon|^{
       \frac{2\beta_0}{\beta_0-k_0-2}
       }}{v_\varepsilon^{
       \frac{2\beta_0}{\beta_0-k_0-2}
      -1}}\nonumber\\
        &\leq  c_6\bigg(\int_{\Omega}w_\varepsilon^{\frac{\beta_0}{\beta_0-k_0-2}}v_\varepsilon+\int_\Omega v_\varepsilon\bigg)\nonumber\\
        &\leq c_6 \|v_0\|_{L^{\infty}(\Omega)}e^{-\kappa t}\big(\int_\Omega w_\varepsilon^{\frac{\beta_0}{\beta_0-k_0-2}}+|\Omega|\big)\\
        &\leq c_7e^{-\kappa t}\|\nabla w_\varepsilon\|_{L^{\frac{11}{4}}(\Omega)}^{\frac{\beta_0 a_0}{\beta_0-k_0-2}}
        \|w_\varepsilon\|_{L^{\frac{5}{2}}(\Omega)}
        ^{\frac{\beta_0(1-a_0)}{\beta_0-k_0-2}}
        +c_7e^{-\kappa t}\|w_\varepsilon\|_{L^2(\Omega)}^{\frac{\beta_0}{\beta_0-k_0+2}}+c_7e^{-\kappa t}\nonumber\\
            &\leq c_8e^{-\kappa t}(1+M^{\frac{2}{3}})^{\frac{\beta_0 a_0}{\beta_0-k_0-2}}+c_8e^{-\kappa t}\nonumber\\
            &\leq c_9 (1+M^{\iota_0})e^{-\kappa t},\nonumber
    \end{align}
    here $q>2$ is warranted by $\beta_0<\beta_2(k_0)<2k_0+4$.
Therefore, combining \eqref{3.11}--\eqref{3.12}, we have
\begin{align}
&~\frac{d}{dt}\bigg(\frac{8\gamma}{3c_1}\int_\Omega u_\varepsilon^{\frac{3}{2}}(\cdot,t)+\frac{
       |\nabla v_\varepsilon|^{
       \frac{2\beta_0}{\beta_0-k_0-2}-2}}{v_\varepsilon^{
       \frac{2\beta_0}{\beta_0-k_0-2}
       -3}}\bigg)+\frac{\gamma}{2c_1}\int_\Omega u_\varepsilon^{\frac{1}{2}}v_\varepsilon|\nabla u_\varepsilon|^2\\
       &\leq \frac{32\gamma}{c_1} \int_\Omega u_\varepsilon^{-\frac{3}{4}}v_\varepsilon|\nabla u_\varepsilon|^2+\frac{c_3(\beta_0)\gamma}{c_1}\int_\Omega v_\varepsilon|\nabla v_\varepsilon|^{\frac{4\beta_0}{2\beta_0-5}}+\frac{c_5\gamma}{c_1}e^{-\kappa t}+c_9 (1+M^{\iota_0})e^{-\kappa t},\nonumber
\end{align}
which, upon integration and using \eqref{3.4} and Lemma \ref{2.6}, shows that
\begin{align}\label{3.14}
       \int_\Omega u_\varepsilon^{\frac{3}{2}}(\cdot,t)&\leq 12\int_0^t\int_{\Omega}u_\varepsilon^{-\frac{3}{4}}v_\varepsilon|\nabla u_\varepsilon|^2\nonumber+c_3(\beta_0)\int_0^t\int_\Omega v_\varepsilon|\nabla v_\varepsilon|^{\frac{4\beta_0}{2\beta_0-5}}+c_5\int_0^t e^{-\kappa t}\\
       &~~+c_9 (1+M^{\iota_0})\int_0^t e^{-\kappa t}+\int_\Omega u_0^{\frac{3}{2}}+\frac{3c_1}{8\gamma}
       \int_\Omega\frac{|\nabla v_0|^{\frac{2\beta_0}{\beta_0-k_0-2}-2}}
       {v_0^{\frac{2\beta_0}{\beta_0-k_0-2}-3}}\\
       &\leq c_{10}(K)(1+M^{\frac{5}{6}})+c_{11}(1+ M^{\iota_0})\int_0^{t}e^{-\kappa s}+c_{12}\int_0^{t}e^{-\kappa s}\nonumber\\
       &~~+\int_\Omega u_0^{\frac{3}{2}}
       +\frac{3c_1}{8\gamma}\int_\Omega\frac{|\nabla v_0|^{\frac{2\beta_0}{\beta_0-k_0-2}-2}}
       {v_0^{\frac{2\beta_0}{\beta_0-k_0-2}-3}}\nonumber\\
       &\leq C(K)(1+M^{\tau})\nonumber
\end{align}
where $\tau:=\max\{\frac{5}{6},\iota_0\}<1$.
%\end{proof}
%\begin{lemma}
 %   Let $\Omega\subset\mathbb{R}^3$ and $K>0$ with the property that \eqref{1.6} holds. There exists some $L=L(m,K)$  such that
 %   \begin{equation}
  %      \int_\Omega u_\varepsilon^{\frac{3}{2}}\leq L~~\mbox{for all}~~t\in(0,\infty).
   % \end{equation}
%\end{lemma}
%\begin{proof}

Now at this position, we take  $M:=(2C(K)+1
)^{\frac{1}{1-\tau}}$, and then have
\begin{equation}\label{3.15}
    \int_\Omega u_\varepsilon^{\frac{3}{2}}(\cdot,t)\leq C(K)(1+ M^{\tau})\leq 2C(K)\cdot M^{\tau}\leq M-\frac 12
\end{equation}
for all $t\in(0,T^*)$. Due to the continuity of $u_\varepsilon$, \eqref{3.15} contradicts with the definition of $T^*$. Hence we have $T^*=\infty$.
\end{proof}

\end{document}
